\section{Discussion}
\label{sec:discussion}

\subsection{Interpreting Reproduction Performance}

\paragraph{The tested object determines the claim.}
The inspected cases show why an executable artifact needs an explicit
target identity. A simulator can implement the expected programming error,
generate realistic responses, and pass its own positive and negative
tests. Those are properties of the simulator. An original firmware
component can provide stronger evidence, yet a direct internal call can
still leave the documented attack entry untested. A useful empirical
account should retain these different scopes instead of forcing them
into a single undifferentiated success label.

This distinction also informs the interpretation of failure. A run that
does not reproduce the documented entry may nevertheless produce a
reusable extracted image, an original-component test, or a verified account
of a missing prerequisite. These artifacts reduce some forms of analyst
work while leaving others unresolved. Their value can be described through
specific verified states without assigning arbitrary exchange rates
between a plan, a download, and a dynamic vulnerability trigger.

\paragraph{Early noncompletion limits what can be inferred later.}
If an agent fails to obtain or execute a target, the record provides little
direct evidence about its ability to construct an exploit once that
target is available. This is a consequence of the observed task path,
not a demonstration that later capability is high. A conditional analysis
of runs that reached execution can help characterize those runs but
selects cases differently for each configuration. It does not remove the
selection problem.

The practical distinction is between deployment performance of the
complete workflow and capability on a specified subtask. Both are useful
questions. They require different experimental conditions to answer
cleanly. Stage-decoupling work makes this explicit by providing missing
prerequisite information and retesting the subsequent
task~\cite{yu2026decoupling}. Our observational analysis will identify
which follow-up comparison would be most informative rather than assume
its result.

\todo{EXP-01}{Optional, if the final claim concerns the effect of
environment readiness: compare CVE-only and verified-ready-target
conditions on the same cases with matched agent setup and budget. Report
the common eligible subset, its selection bias, and negative results.}
Without that experiment, the conclusion should remain that certain
downstream behaviors were not observed after early noncompletion.

\paragraph{Obstacles and responses are complementary explanations.}
An environmental obstacle describes a condition encountered during
execution. A response describes what the agent does in that situation.
The same missing dependency can be followed by a targeted repair,
unchanged repetition, simulation, or termination. Conversely, simulation
can be proposed before any retained execution error. Both dimensions are
needed to avoid attributing every alternative target to a technical
rehosting impossibility.

\todo{DISC-01}{Develop the final cross-RQ interpretation from validated
patterns and counterexamples. Identify which explanation the evidence
supports beyond the individual cases, and retain competing explanations
where unresolved.}
The final interpretation should identify which patterns recur across
cases and which remain isolated examples, and explain how that distinction
changes the practical conclusions.

\subsection{Implications for Agents and Their Evaluation}

\paragraph{Evidence should travel with the result.}
For practical use, a reproduction report should identify the image and
component, how the target ran, which interface received the input, and
what effect was observed. It should disclose modifications that alter
the relevant behavior and separate a simulator's output from original
target execution. This recommendation follows from the different claims
supported by the inspected artifacts; it does not establish that every
existing benchmark would misclassify those artifacts.

A report can further identify the next unresolved prerequisite. For
example, verified component triggering with an unresolved authentication
frontend is more actionable than an unqualified success message. Such
scope information helps a human decide whether to accept a component test,
continue reconstructing the documented interface, or treat the output as
an explanatory demonstration.

\paragraph{Agent improvements should target observed gaps.}
If reviewed trajectories show repeated difficulties interpreting loader
errors or device-state requirements, integration with a specialized
rehosting service is a plausible direction. If a model repeatedly ignores
available tool feedback, another direction is explicit state tracking and
verification before changing strategy. Studies of enterprise penetration
testing and local privilege escalation provide examples of connecting
trajectory observations to architecture or tool
changes~\cite{happe2026enterprisenetworks,huang2025capabilitiesperformance,probst2026localprivesc}.
Our data must establish which of these problems actually recurs in this
setting before recommending a particular implementation.

\todo{EXP-02}{Optional, if claiming that task wording changes simulation
or report alignment: compare the historical instruction with an explicit
original-target requirement that permits honest failure reporting. Keep
the environment and budget matched and measure both progress and claims.}
We do not claim that this change will improve performance or reduce
simulation without testing it. A prompt can alter the task agents believe
they are solving, so any such follow-up should be presented as a new
experimental condition rather than a retrospective correction of the
historical protocol.

\paragraph{Resource value and empirical value can coexist.}
\sys{} names the case collection, protocol, and analysis materials.
The release should let readers inspect each finding and its uncertainty.
External tool comparison remains a separate engineering objective.
